\begin{sectionbox}
	\subsection{Allgemein}
	\begin{tablebox}{ll}
		Volumen & $1l \equalhat 10^{-3}m^3$ \\
		Druck   & $1bar \equalhat 100000Pa$ \\
		Dichte  & $\rho=\frac{m}{V}$ \\
		Reynolds-Zahl & $R_E= \frac{\rho vd}{\eta}$ \\
		& $R_E < R_{E,krit.}$ laminar (z.B. Wasser) \\
		&	$R_E > R_{E,krit.}$ turbulent (z.B. Honig)
	\end{tablebox}
\end{sectionbox}

\begin{sectionbox}
	\subsection{Hydraulische Presse}
	für sich nicht bewegende Flüssigkeiten. Man drückt einen Kolben hinein und ein anderer bewegt sich heraus (z.b. Baggerarm)
	\begin{emphbox}
		$\frac{F_1}{A_1}=\frac{F_2}{A_2}$
	\end{emphbox}
	$\frac{\Delta x_1}{\Delta x_2}=\frac{F_2}{F_1}=\frac{A_2}{A_1}$ mit $\Delta x$: Verschiebung (Hubweg) der Kolben
\end{sectionbox}

\begin{sectionbox}
	\subsection{Kontinuitätsgleichung}
	für sich bewegende Flüssigkeiten (z.B. Wasser im Gartenschlauch wird durch kleine Düsen gedrückt)
	\begin{emphbox}
		$\dot{m}=\rho_1 A_1 \vec{v_1}=\rho_2 A_2 \vec{v_2}$
	\end{emphbox}
	$\dot{V}=\frac{\dot{m}}{\rho}=A\vec{v}$
\end{sectionbox}

\begin{sectionbox}
	\subsection{Bernoulli Gleichung}
	Beschreibt den Gesamtdruck eines Punktes. Dieser setzt sich zusammen aus dem statischen Druck (z.B. atmosphärischer Druck), dem Druck aufgrund einer sich bewegenden Flüssigkeit und dem Druck aufgrund einer Menge an Wasser (z.B. $10m$ unter Wasser)
	\begin{emphbox}
		$p_{\ir ges}=\underbrace{p}_{\text{stat. Druck}}
			+ \underbrace{\frac{1}{2} \rho v^2}_{\text{dyn. D.}:  E_{\ir kin}/V}
			+ \underbrace{\rho gh}_{\text{hydrostat. D.}: E_{\ir pot}/V} = {\small \text{const.}}$
	\end{emphbox}
	$p_{\ir ges,1}\stackrel{!}{=}p_{\ir ges,2}$
	(Energieerhaltung der Hydromechanik)
\end{sectionbox}

\begin{sectionbox}
	\subsection{Gesetz von Hagen-Poiseulle}
	Strömung durch ein Rohr\\
	\begin{emphbox}
		$\dot{V}=\frac{\diff V}{\diff t}=\frac{\pi (p_1-p_2)}{8\eta l}R^4$\\
	\end{emphbox}
	$v(r)=\frac{P_1-P_2}{4\eta l}(R^2-r^2) \qquad \overline{v}=\frac{\dot V}{A}$
\end{sectionbox}

Reynolds-Zahl: TODO: s. Zusammenfassung ZÜ 24.6.

TODO: Rest